\datos{color}{}{}
\section{\textit{ANYmal}}
\textit{ANYmal}, desarrollado por el \textit{Robotic Systems Lab} de la \textit{ETH Zürich} y 
presentado por primera vez en 2016, Figura~\ref{anymal}, es un robot cuadrúpedo avanzado diseñado 
para la inspección autónoma, la exploración en entornos complejos y la intervención en escenarios 
industriales y de rescate. Su diseño modular, su locomoción altamente dinámica y su capacidad para
adaptarse a terrenos irregulares lo han convertido en uno de los robots de patas más influyentes de
la última década. ANYmal integra cámaras estéreo, sensores de profundidad, un escáner láser 3D, 
sensores de fuerza en las articulaciones y un ordenador principal de alto rendimiento encargado de
la percepción, la planificación y la navegación. \\

Sus usos principales incluyen la inspección autónoma en plantas industriales, la exploración de
túneles y minas, la cartografía tridimensional, la investigación en control locomotor y la 
participación en competiciones internacionales como la \textit{DARPA Subterranean Challenge}. 
Gracias a su robustez mecánica y a su capacidad para caminar, trotar, subir escaleras, levantarse
tras caídas y superar obstáculos complejos, ANYmal se ha convertido en una plataforma de referencia 
para la investigación en robótica cuadrúpeda. \\

Al igual que otros robots cuadrúpedos modernos como \textit{Spot} o \textit{Unitree Go1}, 
\textit{ANYmal} implementa una arquitectura de control híbrida. El ordenador central ejecuta los 
módulos de percepción, integración multisensorial, planificación de trayectorias y navegación global,
siguiendo la descomposición vertical clásica. Sin embargo, cada articulación dispone de controladores
locales que procesan directamente la información sensorial y generan respuestas rápidas para mantener
la estabilidad dinámica, absorber impactos o corregir perturbaciones. Los niveles superiores
planifican el movimiento global y los objetivos de navegación, mientras que los niveles inferiores
garantizan la ejecución reactiva y segura del patrón de marcha. Esta combinación de planificación
deliberativa y control reactivo permite a ANYmal adaptarse rápidamente a cambios del entorno,
cumpliendo las características de la arquitectura híbrida descrita en la Sección~2.3. \\

\begin{wrapfigure}{r}{0.34\textwidth}
\centering
\vspace{-10pt}
\includegraphics[width=0.34\textwidth]{Images/anymal.jpg}
\caption{\label{anymal}\textit{ANYmal}.}
\vspace{-10pt}
\end{wrapfigure}

\textit{ANYmal} utiliza un sistema de locomoción cuadrúpedo con doce grados de libertad, 
accionados mediante motores eléctricos de alta potencia y sensores de torque. La combinación
de control de impedancia, estimación del estado corporal y planificación dinámica permite al
robot caminar, correr, girar con precisión, levantarse tras caídas y desplazarse por terrenos
irregulares. Su locomoción reactiva y su capacidad para mantener la estabilidad ante perturbaciones
lo convierten en un ejemplo representativo de las arquitecturas híbridas modernas aplicadas a robots
cuadrúpedos de alto rendimiento. \\

\subsection*{Referencias}
\cite{ANYmal2016} \fullcite{ANYmal2016}.\\
\cite{ANYmalSubT} \fullcite{ANYmalSubT}.